problem:  
Let \( G \) be a compact, connected, simple Lie group equipped with a one-parameter family of **left-invariant naturally reductive** Riemannian metrics defined by  
\[
g_t(X,Y) = \langle X_{\mathfrak{k}}, Y_{\mathfrak{k}} \rangle + t \, \langle X_{\mathfrak{p}}, Y_{\mathfrak{p}} \rangle , \quad t > 0,
\]
where \(\mathfrak{g} = \mathfrak{k} \oplus \mathfrak{p}\) is an \( \mathrm{Ad}(K)\)-invariant decomposition for some proper subgroup \(K \subset G\) and \(\langle \cdot,\cdot\rangle\) is the negative of the Killing form of \(\mathfrak{g}\). (Such families are known to remain naturally reductive with respect to \(G \times K\).)

For each \(t\), denote by \(J_t(X) = R_t(X, \cdot)X\) the Jacobi operator of the metric \(g_t\), where \(R_t\) is its curvature tensor.  

**Conjectural problem:**   
Assume that, for every unit vector \(X \in \mathfrak{g}\), the **spectrum of the Jacobi operator** \(J_t(X)\) (counting multiplicities) is *independent of \(t\)*.  
Must it follow that \((G,g_1)\) is locally symmetric (and hence isometric to \(G\) with its bi-invariant metric)?  

Equivalently, does invariance of curvature operator spectra under all left-invariant naturally reductive deformations characterize bi-invariant metrics among compact simple Lie groups?

Why is it a "good" problem:  
This problem embodies a refined **spectral rigidity conjecture** in the context of naturally reductive homogeneous geometries. It is formulated within a concrete, nontrivial class (compact simple Lie groups) where naturally reductive deformations are known and curvature can be computed explicitly, ensuring the problem is well-posed and non-vacuous.

It is “good” because:

1. **Conceptually simple and natural:**  
   The question can be stated succinctly—“Does spectral invariance imply symmetry?”—yet it cuts to the core of how curvature determines geometry in homogeneous spaces.

2. **Deep structural interaction:**  
   Resolving the problem forces a synthesis of **Lie algebraic curvature computations**, **representation-theoretic constraints** on isotropy modules, and **spectral geometry**. One must understand how deformation of structure constants affects the entire spectrum of the Jacobi operator.

3. **Extension of known frameworks:**  
   It generalizes the **Osserman problem**, which characterizes rank-one symmetric spaces from directional spectral invariance, by replacing “directional independence” with “parametric invariance.” It thereby introduces a new form of spectral rigidity tied not to tangent direction but to global geometric deformation.

4. **Potential payoff:**  
   A positive resolution would yield a **spectral characterization of bi-invariant metrics** and connect curvature invariants to the intrinsic symmetries of Lie groups. A counterexample would expose unknown flexibility in the spectral geometry of naturally reductive metrics.

5. **Bridging disciplines:**  
   The question lies at the junction of **spectral geometry**, **homogeneous Riemannian geometry**, and **Lie theory**. Any partial progress would require advancing understanding in all three simultaneously.

Hence, the problem is specific enough to be feasible, profound enough to remain open, and simple enough in formulation to exemplify the “narrow entrance, profound depth” character of a truly “good” problem.